%%%%%%%%%%%%%%%%%%%%%%%%%%%%%%%%%%%%%%%%%%%%%%%%%%%%%%%%%%%%%%%%%%%%%%%%%%%%%%%%
%%%%%%%%%%%%%%%%%%%%%%%%%%%%%%%%%%%%%%%%%%%%%%%%%%%%%%%%%%%%%%%%%%%%%%%%%%%%%%%%
% ESTADO DEL ARTE %
% Tecnmolofgias q usas para el TFG (todo lo q no has hecho tu : Lenmguajes de programacion, herramientas, bibliotecas etc)%
%%%%%%%%%%%%%%%%%%%%%%%%%%%%%%%%%%%%%%%%%%%%%%%%%%%%%%%%%%%%%%%%%%%%%%%%%%%%%%%%

\cleardoublepage
\chapter{Estado del arte}
\label{chap:estado}

En este capitulo describiré las tecnologías, herramientas y conceptos mas importantes que hicieron posible la creacion de WebBuilder, es decir, todo aquello que no ha sido creado en el marco de este trabajo pero que son fundamentales.

Cubriremos desde el framework web y la base de datos que sustentan la aplicacion, hasta los modelos de lenguaje y las herramientas de automatizacion e infraestructura que hacen posible el funcionamiento del proyecto.

El capítulo se organiza en cuatro bloques temáticos. En primer lugar se presentan las tecnologías de backend y base de datos que conforman el núcleo de la aplicación. A continuación se introduce el campo de la inteligencia artificial generativa y los modelos de lenguaje, pieza central del proyecto, junto con los conceptos de ingeniería de prompts y los proveedores de API utilizados. El tercer bloque cubre n8n, la herramienta de automatización que orquesta los flujos de notificación y despliegue. Por último, se describen Docker y Tailwind CSS, tecnologías de infraestructura y presentación empleadas en los proyectos web generados por la plataforma.



\section{Backend} 
\label{sec:backend}

El backend está construido con diferentes tecnologías de desarrollo web, en esta sección se describen las herramientas que forman la capa de servidor y persistencia de datos de la aplicación.


\subsection{Python}
\label{subsec:python}

Python es un lenguaje de programación interpretado, de alto nivel y útil para diferentes tipos de tareas. Creado por Guido van Rossum a principios de los años 90~\cite{python_history}, Python ha evolucionado hasta convertirse en uno de los lenguajes más utilizados del mundo, con un 57,9\% de uso entre los desarrolladores en 2025~\cite{stackoverflow_survey_2025}, especialmente en el ámbito
del desarrollo web, el análisis de datos y la inteligencia artificial. Siendo uno de los lenguajes de programación más versátiles.

La sintaxis clara permite escribir programas concisos para combinar lógica del servidor, procesamiento de APIs y generación dinámica de código. Python cuenta con un ecosistema de bibliotecas muy completo, distribuido a través de PyPI (\emph{Python Package Index}), del que WebBuilder hace uso directo: \texttt{requests} para la descarga de datos desde APIs externas, \texttt{xmltodict} y \texttt{defusedxml} para el parseo seguro de XML y \texttt{zipfile} para el empaquetado de los proyectos generados.


En WebBuilder, Python actúa como columna vertebral de la lógica del servidor, gestionando la ingestión y validación de datos, la construcción y envío de prompts al LLM y la generación del código de los proyectos. Los detalles técnicos de cada función se describen en el Capítulo~\ref{chap:diseno}.

La versión utilizada es Python~3, aprovechando características modernas del lenguaje
como las \emph{f-strings} (cadenas de texto con expresiones Python embebidas, usadas en la construcción de prompts para el LLM).

\subsection{Django}
\label{subsec:django}

Django es un framework de desarrollo web de código abierto para Python, que promueve 
el desarrollo rápido y escalable, así como un diseño limpio~\cite{django_docs}. 
Fue publicado en 2005 y sigue el principio \emph{``batteries included''}: incluye 
de serie todo lo necesario para construir aplicaciones web completas, desde el sistema 
de enrutado de URLs hasta el ORM (\emph{Object-Relational Mapper}, capa de software que permite trabajar con la base de datos usando objetos Python en lugar de escribir SQL), el motor de 
plantillas, el sistema de autenticación y el panel de administración.

Permite configurar todo lo necesario, Django incluye un sistema de gestión de base de datos SQLite integrado que permite definir las tablas y sus campos directamente en Python, sin necesidad de escribir SQL. Cada tabla de la base de datos se representa como una clase, y Django se encarga de traducir automáticamente las operaciones sobre esos objetos (consultas, inserciones, actualizaciones) a las instrucciones SQL correspondientes. Esto facilita el desarrollo inicial, y también la evolución del modelo a través de migraciones automáticas, que registran cada cambio de forma controlada, siendo capaz de saber que se cambio o modifico en cada momento, teniendo más control sobre el desarrollo.

La lógica de vistas y de URLs permite organizar la lógica de forma clara y separada, con funciones clave que responden peticiones HTTP, separando el procesamiento de datos de la presentación. El motor de plantillas permite generar HTML dinámico con una sintaxis sencilla, con soporte para herencia de plantillas, lo que favorece la reutilización de componentes visuales.

En WebBuilder, Django gestiona el ciclo de vida completo de cada análisis: desde la recepción y validación de la URL hasta la persistencia de resultados, la invocación del módulo LLM y el servicio de las vistas del asistente. Su funcionamiento detallado se describe en el Capítulo~\ref{chap:diseno}.

La versión empleada es Django~5.1, que introduce mejoras en el ORM y en el sistema de migraciones respecto a versiones anteriores.


\subsection{Base de datos: SQLite y PostgreSQL}
\label{subsec:base-datos}

SQLite es un sistema de base de datos ligero que almacena toda la información en un único fichero local, sin necesidad de instalar ni configurar un servidor externo~\cite{sqlite_docs}. Esta simplicidad lo convierte en la opción por defecto de Django para entornos de desarrollo, aunque presenta limitaciones en escrituras concurrentes al utilizar bloqueos sobre el fichero de datos.

PostgreSQL es un sistema gestor de bases de datos relacional de código abierto, robusto y ampliamente utilizado en entornos de producción~\cite{postgresql_docs}. Ofrece soporte completo para transacciones ACID, concurrencia real mediante MVCC (\emph{Multi-Version Concurrency Control}) y un rico conjunto de tipos de datos, lo que lo convierte en una opción adecuada cuando se requiere mayor robustez y escalabilidad.

En WebBuilder se utilizaron ambas tecnologías de forma progresiva: SQLite durante las fases iniciales de desarrollo por su sencillez, y PostgreSQL en las fases finales para soportar un entorno de producción más exigente. El detalle de la migración de SQLite a PostgreSQL y los motivos técnicos que la justificaron se describen en la Sección~\ref{subsec:migracion-postgresql}.


\section{Inteligencia artificial y modelos de lenguaje}
\label{sec:ia}

La inteligencia artificial generativa constituye una parte importante del motor del proyecto. En esta sección se introducen los conceptos fundamentales sobre los modelos de lenguaje, el prompting engineering y los proveedores de API utilizados para acceder a los diferentes modelos.


\subsection{Modelos de lenguaje y IA generativa}
\label{subsec:llms}

Los modelos de lenguaje de gran tamaño (LLM, \emph{Large Language Models}) son sistemas de inteligencia artificial entrenados sobre grandes volúmenes de datos. Su arquitectura se basa en redes neuronales de tipo \emph{Transformer}~\cite{attention_is_all_you_need}, introducidas en 2017, que sustituyeron el procesamiento secuencial del texto por un mecanismo de atención capaz de analizar todas las palabras simultáneamente y capturar relaciones de largo alcance entre ellas.

A diferencia de los sistemas tradicionales de procesamiento de lenguaje natural, los LLM no se programan con reglas explícitas, sino que aprenden a predecir la siguiente unidad de texto, denominada token, a partir de las anteriores. Así es como se da lugar a modelos capaces de mantener conversaciones, redactar textos, traducir idiomas o como en este caso, generar código funcional a partir de instrucciones en lenguaje natural.
Cada modelo tiene una ventana de contexto llamada \textbf{context window}, que limita el número máximo de tokens que puede procesar simultáneamente, incluyendo tanto el prompt de entrada como la respuesta generada. Superar este límite implica que el modelo pierde información de los mensajes más antiguos. Este límite es muy importante, ya que los prompts de generación incluyen el esquema de datos, ejemplos del dataset e instrucciones de formato, lo que puede consumir varios miles de tokens por llamada, como se detalla en el capítulo~\ref{chap:diseno}.

Un fenómeno frecuente en los LLM son las alucinaciones, estas ocurren cuando el modelo genera contenido aparentemente valido pero en realidad es incorrecto o directamente inexistente~\cite{hallucination_survey}. A diferencia de un error de compilación, una alucinación no produce un fallo sino una salida que parece válida pero que contiene datos inventados, nombres de campos que no existen en el dataset o código erróneo. Las alucinaciones son el principal riesgo del sistema, ya que un campo inventado en el modelo Django generado provoca errores silenciosos en tiempo de ejecución. El diseño de los mecanismos de protección descritos en la Sección~\ref{subsec:proteccion-anti-alucionaciones} responden directamente a este problema.


La irrupción mediática de modelos como ChatGPT en 2022~\cite{chatgpt_launch} marcó 
el inicio de una etapa de adopción masiva de la IA generativa. Desde entonces han 
aparecido numerosos modelos, tanto comerciales (GPT-4, Claude, Gemini) como de 
pesos abiertos (LLaMA, Mistral, DeepSeek), accesibles a través de APIs públicas 
o ejecutables localmente. Esta diversidad permite a WebBuilder elegir el modelo 
más adecuado para cada tarea en función de criterios como coste, latencia o 
calidad de respuesta.

En este trabajo, los LLM no se utilizan como simples chatbots, sino como componentes programables dentro de un flujo automatizado: el modelo recibe datos estructurados, devuelve análisis razonados y produce código fuente que se integra directamente en el proyecto generado.


\subsection{Ingeniería de prompts}
\label{subsec:prompts}

La ingeniería de prompts (\emph{prompt engineering}) es la disciplina que estudia cómo formular las instrucciones enviadas a un modelo de lenguaje para obtener respuestas precisas y acorde a una necesidad, intentando a su vez reducir el numero de alucinaciones. A diferencia de la programación tradicional, donde el comportamiento del sistema se especifica mediante código, en los LLM el comportamiento se controla a través del lenguaje natural del propio prompt.

Para poder diseñar un buen prompt debemos seguir y combinar varios elementos: 

\begin{itemize}
    \item \textbf{Rol del modelo:} una descripción clara de la identidad o 
    función que debe asumir el modelo durante la tarea.
    \item \textbf{Instrucciones detalladas:} una especificación precisa de 
    lo que se espera que el modelo haga y cómo debe hacerlo.
    \item \textbf{Formato de respuesta:} indicaciones explícitas sobre la 
    estructura que debe tener la salida, como JSON, código Python o texto 
    plano.
    \item \textbf{Restricciones:} limitaciones explícitas que acotan el 
    comportamiento del modelo y reducen la probabilidad de respuestas 
    incorrectas o fuera de contexto.
    \item \textbf{Ejemplos de referencia:} muestras del estilo de salida 
    esperado que guían al modelo hacia el resultado deseado. Esta técnica, 
    conocida como \emph{few-shot prompting}, ha demostrado mejorar 
    notablemente la calidad de las respuestas en tareas complejas o con 
    formatos estrictos~\cite{openai_prompt_engineering}.
\end{itemize}


La calidad del resultado generado por un LLM depende, en \textbf{gran medida}, de la precisión del prompt. Un prompt ambiguo puede provocar respuestas inconsistentes, alucinaciones o salidas que no respetan el formato esperado. Por ello, el diseño, la iteración y el refinamiento de los prompts constituye una parte central del trabajo de desarrollo en sistemas basados en LLM, comparable en importancia a la propia escritura del código que los invoca.

La calidad del resultado no depende únicamente de la precisión del prompt, sino también de la capacidad del modelo utilizado. Un modelo con poco entrenamiento 
o potencia insuficiente producirá respuestas pobres independientemente de la calidad del prompt. Ambos factores, diseño del prompt y elección del modelo, son por tanto determinantes en el resultado final.

En el contexto de WebBuilder, la ingeniería de prompts juega un papel crítico en dos fases: el análisis de los datos de la API y la generación del código Django del proyecto. Cada uno de estos prompts ha sido diseñado tras innumerables pruebas para producir resultados coherentes y de calidad profesional, como se detallará en el Capítulo~\ref{chap:diseno}. Además, durante el proyecto se fueron guardando los datos de cada generación para realizar un estudio exhaustivo de qué funciona mejor para cada caso, cuyos resultados se recogen en el Capítulo~\ref{chap:experimentos}.


\subsection{Proveedores de modelos}
\label{subsec:proveedores}

Para acceder a un LLM y poder usarlo, en la mayoría de casos es necesario recurrir a un proveedor de API que disponga del modelo y ofrezca su uso. Cada vez existen mas proveedores en el mercado, cada uno con su propio catalogo. La gran mayoría de estos proveedores siguen el estándar de OpenAI descrito a continuación, lo que permite conectar de forma sencilla con WebBuilder. Cabe aclarar que los modelos gratuitos de los proveedores mencionados a continuación realizan bien su trabajo, pero nunca de forma perfecta. Las alucinaciones a día de hoy siguen existiendo sobre todo en los modelos gratuitos.

\textbf{Groq}~\cite{groq_docs} es un proveedor especializado en inferencia de baja latencia. Utiliza hardware propio (LPU, \emph{Language Processing Units}) optimizado, lo que le permite ofrecer velocidades de generación muy superiores a las de proveedores tradicionales basados en GPU. Groq aloja modelos de pesos abiertos como LLaMA, y dispone de un nivel de uso gratuito con límites de peticiones por minuto y por día. Su baja latencia hace que sea la opción preferente, donde la generación de templates es rapidísima.

\textbf{OpenRouter}~\cite{openrouter_docs} es un proveedor que da acceso a decenas de modelos de distintos proveedores a través de una única API. Lo mejor de este proveedor es la sencillez para la conexión: una sola clave de API permite consultar modelos de OpenAI, Anthropic, Google, Meta o cualquier otro proveedor integrado en la plataforma. Esto facilita las pruebas comparativas entre modelos y reduce el coste de cambiar de proveedor en caso necesario.

Otros proveedores como \textbf{NVIDIA NIM}~\cite{nvidia_nim_docs}, también compatibles con este estándar, amplían el catálogo con modelos especializados en código y razonamiento.

Los tres proveedores descritos comparten el denominado formato \textit{OpenAI-compatible}, un estándar surgido a partir de la API pública de OpenAI~\cite{openai_api_docs}. Este estándar define un único endpoint \texttt{/chat/completions} que acepta una lista de mensajes con roles diferenciados (\texttt{system}, \texttt{user} y \texttt{assistant}) y devuelve la respuesta del modelo en un formato JSON uniforme. Su adopción masiva por parte de los principales proveedores del mercado lo ha convertido en el protocolo común de la industria de modelos de lenguaje.

Esto es algo muy importante para el proyecto: al construir el cliente LLM sobre este protocolo común, el sistema puede cambiar de proveedor o de modelo modificando únicamente tres parámetros de configuración (la URL base, la clave de API y el identificador del modelo), sin necesidad de modificar ninguna línea de la lógica interna. Esta decisión de diseño, detallada en la Sección~\ref{subsec:conexion-llm}, es la que hace posible que WebBuilder soporte múltiples proveedores de forma nativa y que cualquier usuario pueda conectar su propio modelo personalizado.
\section{Automatización e infraestructura}
\label{sec:auto}

Como se comento anteriorme, WebBuilder no solo se apoya en el backend de Django y Pyhton, tambien se usan herramientas extenas para la automatizacion de tareas y el despliegue de los proyectos. En esta sección se describen las herramientas de n8n, encargado de orquestar los flujos asíncronos de la plataforma, y de Docker, que proporciona el entorno aislado en el que se ejecutan los proyectos generados.


\subsection{n8n}
\label{subsec:n8n}

n8n es una plataforma de automatización de flujos de trabajo de código abierto que permite conectar aplicaciones, APIs y servicios mediante una interfaz visual basada en nodos. Cada flujo de trabajo (\emph{workflow}) está compuesto por una secuencia de nodos quese ejecutan en cadena, donde cada nodo representa una accion, explicando su funcion dentro del contenido del mismo nodo, hay varias funciones: 

\begin{itemize}
    \item \textbf{Disparadores:} eventos que inician el flujo, como una 
    petición HTTP entrante (webhook) o un temporizador programado (cron).
    \item \textbf{Peticiones HTTP:} llamadas a APIs externas o internas para 
    intercambiar datos entre servicios.
    \item \textbf{Transformación de datos:} procesamiento y manipulación de 
    la información entre nodos mediante lógica personalizada en JavaScript.
    \item \textbf{Envío de correos:} notificaciones automáticas por email 
    a usuarios o administradores ante determinados eventos del sistema.
    \item \textbf{Ejecución de scripts:} bloques de código personalizado 
    que permiten implementar lógica compleja directamente dentro del flujo.
    \item \textbf{Consultas a base de datos:} lectura o escritura de datos 
    en sistemas de persistencia externos integrados en el workflow.
\end{itemize}

Los flujos pueden pueden iniciarse a partir de webhooks (peticiones HTTP entrantes), disparadores temporales (\emph{cron}) o eventos de aplicaciones conectadas, y muchos de los nodos implementados en este proyecto son de logica de programacion. Estos nodos permite ejecutar fragmentos de codigo JavaScript personalizado, implementando logicas completas como la limpieza y despliegue de contenedores Docker.

En WebBuilder, n8n actúa como una capa desacoplada que gestiona todo lo 
que no forma parte del ciclo síncrono de petición-respuesta de Django:

n8n actúa como una capa desacoplada que gestiona todo lo que no forma parte del ciclo síncrono de petición-respuesta de Django: 
el envío de correos automáticos al registrarse o iniciar sesión, la notificación al finalizar una generación, la monitorización diaria del estado del sistema, el apagado automático de contenedores inactivos y, sobre todo, el despliegue de los proyectos generados en contenedores Docker. Esta arquitectura permite añadir nuevas automatizaciones sin modificar el código del backend principal.

A diferencia de plataformas similares, n8n puede autoalojarse, lo que permite ejecutarlo en infraestructura propia y mantener el control completo sobre los datos y las credenciales. Como ocurre en este caso que esta desplegado dentro de un contenedor en local, pudiendo acceder y tener el control total de la herramienta


\subsection{Docker}
\label{subsec:docker}

Dcoker es una herramienta que permite empaquetar aplicaciones junto con sus dependencias en unidades aisladas llamadas contenedores.
A diferencia de la virtualización tradicional, los contenedores comparten el núcleo del sistema operativo anfitrión, lo que los hace mucho más ligeros y rápidos de arrancar que una máquina virtual convencional.

Cada contenedor se construye a partir de una imagen, un artefacto inmutable que contiene el sistema de ficheros, el software instalado y la configuración necesaria para ejecutar la aplicación. Las imágenes se definen mediante un fichero de texto llamado \texttt{Dockerfile}, que 
describe paso a paso cómo construir el entorno: la imagen base, los paquetes a instalar, el código a copiar y el comando que se ejecuta al arrancar el contenedor.

Esta capacidad de empaquetado garantiza que la aplicación se ejecute de forma idéntica en cualquier entorno: el equipo de desarrollo, un servidor de pruebas o un entorno de producción. Esto resuelve uno de los problemas clásicos del despliegue, ejecutandose correctamente en algunas maquinas pero en otras no, debido a las diferencias entre los entornos en los que se desarrolla y se ejecuta una aplicación.

En WebBuilder, Docker cumple dos funciones diferenciadas. Por un lado, aloja la propia herramienta de n8n utilizada por la plataforma de forma local. Por otro, y más importante, cada proyecto Django generado por WebBuilder incluye automáticamente su propio \texttt{Dockerfile}, lo que permite desplegarlo en un contenedor aislado sin necesidad de instalar dependencias en el sistema anfitrión. Un flujo especifico de n8n (webbuilder-deploy), construye la imagen, lanza el contenedor y le 
asigna un puerto disponible de forma totalmente automática.

\section{Frontend}
\label{sec:frontend}

WebBuilder tiene dos capas de frontend diferenciadas: la interfaz de la propia 
plataforma, construida con tecnologías web estándar, y los proyectos que genera.

En este primer caso, la interfaz de WebBuilder está construida con HTML, CSS y 
JavaScript, las tres tecnologías fundamentales de la web, estandarizadas por el 
W3C (\emph{World Wide Web Consortium})~\cite{w3c}. HTML define la estructura de 
cada página, CSS controla su presentación visual y JavaScript gestiona la 
interactividad en el lado del cliente.

\subsection{HTML, CSS y JavaScript}
\label{subsec:htmlcssjs}

\textbf{HTML}~\cite{mdn_html} (\emph{HyperText Markup Language}) es el lenguaje de marcado que define la estructura y el contenido de las páginas web. En WebBuilder, todas las plantillas Django están escritas en HTML, organizando los elementos visuales del asistente, los formularios, las tablas de historial y los paneles de información.

\textbf{CSS}~\cite{mdn_css} (\emph{Cascading Style Sheets}) es el lenguaje de estilos que controla la presentación visual del HTML: colores, tipografías, espaciados,  disposición de los elementos y diseño responsive. La interfaz de WebBuilder utiliza CSS para definir tanto el tema visual general de la plataforma como los modos claro y oscuro implementados durante el desarrollo.

\textbf{JavaScript}~\cite{mdn_javascript} es el lenguaje de programación que añade interactividad en el lado del cliente, permitiendo modificar el contenido de la página sin necesidad de recargarla. En WebBuilder se emplea en varias vistas clave: la previsualización dinámica del plan propuesto por el LLM, el editor de código integrado para inspeccionar los ficheros generados y la actualización del estado del asistente en cada paso del flujo.



\subsection{Tailwind CSS}
\label{subsec:tailwind}

Tailwind CSS~\cite{tailwind_docs} es un framework de hojas de estilo basado en el enfoque \emph{utility-first}. A diferencia de frameworks tradicionales como  Bootstrap, que ofrecen componentes predefinidos (botones, tarjetas, formularios) listos para usar, Tailwind proporciona un conjunto extenso de clases  que se combinan directamente en el HTML para construir cualquier diseño desde cero.

Este enfoque ofrece varias ventajas prácticas. En primer lugar, evita el problema de los estilos huérfanos y la acumulación de CSS no utilizado, ya que cada clase aplicada está explícitamente presente en el marcado. En segundo lugar, permite un control fino sobre el diseño sin necesidad de sobrescribir estilos predefinidos, lo que reduce la complejidad del CSS final. Por último, su sistema de diseño coherente con escalas predefinidas para espaciados, colores, tipografías y puntos de ruptura responsivos favorece la consistencia visual entre componentes.

Tailwind se ha convertido en una de las opciones más populares para proyectos 
web modernos~\cite{stackoverflow_survey_2025}, y su naturaleza declarativa encaja 
especialmente bien con la generación automática de interfaces por parte de modelos 
de lenguaje, ya que los LLM tienden a producir HTML con clases Tailwind de forma 
natural gracias a la abundancia de ejemplos presentes en sus datos de entrenamiento.

En WebBuilder, todos los proyectos Django generados incluyen Tailwind CSS como 
sistema de estilos por defecto, garantizando una apariencia visual coherente 
y profesional sin necesidad de mantener hojas de estilo personalizadas para 
cada proyecto generado.

